\emph{Översikt}
Hittills har programmen räknat på värden som stod skrivna i koden: ett
namn, ett årtal, en längd.  Ett sådant program gör exakt samma sak varje
gång det körs.  Den här modulen gör programmen \emph{interaktiva}: de
frågar den som kör dem.  Vi använder \mintinline{python}{input} för att
läsa in text, och upptäcker att det som kommer tillbaka alltid är en
\emph{sträng} --- även när användaren skriver siffror.  Därför behöver vi
\emph{typkonvertering}, och därför behöver vi veta vad datatypen gör med
en operation: \mintinline{python}|"3" + "4"| och
\mintinline{python}{3 + 4} ser lika ut men ger olika svar.  Vi går
igenom de vanligaste operationerna på tal och strängar, och ett par
strängmetoder som vi får användning för när vi ska granska det användaren
skrev.  Sist utformar vi frågorna och utskrifterna så att den som kör
programmet förstår vad som förväntas.

\emph{Lärandemål}
Efter denna modul ska studenten kunna:
\begin{restatable}{lo}{InputLORead}\label{InputLORead}%
  Läsa in text från användaren med \mintinline{python}{input} och
  förklara att resultatet alltid är en sträng, oavsett vad användaren
  skrev.
\end{restatable}% Lo08
\begin{restatable}{lo}{InputLOConvert}\label{InputLOConvert}%
  Omvandla en inläst sträng till \mintinline{python}{int} eller
  \mintinline{python}{float} med typkonvertering, och avgöra var i
  programmet omvandlingen hör hemma.
\end{restatable}% Lo08, Lo03
\begin{restatable}{lo}{InputLOTypes}\label{InputLOTypes}%
  Använda de vanligaste operationerna och metoderna på tal och strängar,
  och förutsäga vilken typ och vilket värde resultatet får.
\end{restatable}% Lo03
\begin{restatable}{lo}{InputLOPrompt}\label{InputLOPrompt}%
  Utforma frågor och utskrifter så att den som kör programmet förstår
  vad som förväntas och vad svaret betydde.
\end{restatable}% Lo05

\emph{Förkunskaper}
Modulen bygger på föreläsningen \emph{Funktioner} från veckan innan:
variabler, datatyperna \mintinline{python}{int}, \mintinline{python}{float}
och \mintinline{python}{str}, utskrifter med \mintinline{python}{print}
och egna funktioner med parametrar och returvärden.

\emph{Läsning}
Kompletterande material finns i kursens videogenomgång \emph{Inmatning
och felhantering}.  Vad programmet ska göra när användaren skriver något
oväntat --- särfall och \mintinline{python}{try}/\mintinline{python}{except}
--- tas upp i modulen om felhantering.
